% Источники: 09_Архитектура_и_API/DOC-07_Архитектура.md
% Генератор: docs/build_pdf.py
% Шаблон печатного представления. Содержимое подставляет docs/build_pdf.py.
\documentclass[12pt,a4paper]{article}
\usepackage{iftex}
\ifPDFTeX
  \usepackage[utf8]{inputenc}
  \usepackage[T2A]{fontenc}
  \usepackage{newtxtext,newtxmath}
\else
  \usepackage{fontspec}
  \IfFontExistsTF{Times New Roman}
    {\setmainfont{Times New Roman}}
    {\setmainfont{Libertinus Serif}}
  \IfFontExistsTF{Arial}{\setsansfont{Arial}}{\setsansfont{Latin Modern Sans}}
  \IfFontExistsTF{Menlo}{\setmonofont{Menlo}[Scale=MatchLowercase]}
    {\setmonofont{Latin Modern Mono}[Scale=MatchLowercase]}
\fi
\usepackage[russian]{babel}
\usepackage{geometry}
\geometry{left=25mm,right=15mm,top=20mm,bottom=20mm}
\usepackage{setspace}
\onehalfspacing
\usepackage{indentfirst}
\setlength{\parindent}{1.25cm}
\setlength{\parskip}{2pt}
\setlength{\emergencystretch}{3em}
\usepackage{graphicx,booktabs,longtable,array,tabularx,calc,etoolbox}
\usepackage{enumitem}
\setlist[itemize]{leftmargin=1.25cm,itemsep=2pt,topsep=4pt}
\setlist[enumerate]{leftmargin=1.25cm,itemsep=2pt,topsep=4pt}
\newcommand{\tightlist}{\setlength{\itemsep}{0pt}\setlength{\parskip}{0pt}}
\AtBeginEnvironment{longtable}{\small\singlespacing}
\setlength{\LTpre}{8pt}
\setlength{\LTpost}{8pt}
\setlength{\tabcolsep}{5pt}
\setlength{\LTleft}{0pt}
\setlength{\LTright}{0pt}
\usepackage{xurl}
\usepackage[unicode,colorlinks=true,linkcolor=black,urlcolor=blue]{hyperref}
\usepackage{bookmark}
\hypersetup{pdftitle={Архитектура ПО},pdfauthor={Георгиев Д.Д.},pdfsubject={Криптотренажёр mephi-12}}
\urlstyle{same}
\pagestyle{plain}
\setcounter{secnumdepth}{3}
\setcounter{tocdepth}{2}

\begin{document}
\hypersetup{pageanchor=false}
\begin{titlepage}
  \thispagestyle{empty}
  \centering
  {\large НИЯУ МИФИ\par}
  \vspace*{3.5cm}
  {\LARGE\bfseries Архитектура ПО\par}
  \vspace{0.5cm}
  {\large DOC-07\par}
  \vspace{0.6cm}
  {\large Платформа для изучения и практики курса по криптографии\par}
  {\large Криптотренажёр mephi-12\par}
  \vspace{0.7cm}
  {\large Версия 0.5\par}
  {\large Редакция от 07.11.2026\par}
  \vspace{1cm}
  \begin{tabularx}{\textwidth}{@{}p{3.4cm}Xp{4cm}@{}}
    \toprule
    Роль & ФИО & Подпись / Дата \\
    \midrule
    Разработал & Георгиев Д.Д. & \\
    \addlinespace[6pt]
    Согласовано & & \\
    \addlinespace[6pt]
    Утверждено & & \\
    \bottomrule
  \end{tabularx}
  \vfill
  {\large Москва, 2026\par}
\end{titlepage}
\hypersetup{pageanchor=true}

\section*{Лист регистрации изменений}
\addcontentsline{toc}{section}{Лист регистрации изменений}
\begingroup\small\singlespacing
\noindent
\begin{tabularx}{\textwidth}{@{}p{1.2cm}p{2cm}p{2.5cm}Xp{1.8cm}@{}}
\toprule
Версия & Дата & Автор & Содержание изменения & Утвердил \\
\midrule
0.5 & 07.11.2026 & Георгиев Д.Д. & Подготовлен проект архитектуры: слои, компоненты и связи, развёртывание, соответствие требованиям. & \\
\bottomrule
\end{tabularx}
\endgroup
\clearpage
\begingroup\small\singlespacing
\tableofcontents
\endgroup
\clearpage

\hypertarget{doc-07-архитектура-по}{}

{\def\LTcaptype{none} % do not increment counter
\begin{longtable}[]{@{}
  >{\raggedright\arraybackslash}p{(\linewidth - 2\tabcolsep) * \real{0.3200}}
  >{\raggedright\arraybackslash}p{(\linewidth - 2\tabcolsep) * \real{0.6800}}@{}}
\toprule\noalign{}
\begin{minipage}[b]{\linewidth}\raggedright
Поле
\end{minipage} & \begin{minipage}[b]{\linewidth}\raggedright
Значение
\end{minipage} \\
\midrule\noalign{}
\endhead
\bottomrule\noalign{}
\endlastfoot
Версия & 0.5 \\
Дата & 07.11.2026 \\
Ответственный & Георгиев Д.Д. \\
Основания & \href{../07_\%D0\%9F\%D1\%80\%D0\%BE\%D0\%B3\%D1\%80\%D0\%B0\%D0\%BC\%D0\%BC\%D0\%BD\%D1\%8B\%D0\%B5_\%D1\%82\%D1\%80\%D0\%B5\%D0\%B1\%D0\%BE\%D0\%B2\%D0\%B0\%D0\%BD\%D0\%B8\%D1\%8F/DOC-06_\%D0\%9F\%D1\%80\%D0\%BE\%D0\%B3\%D1\%80\%D0\%B0\%D0\%BC\%D0\%BC\%D0\%BD\%D1\%8B\%D0\%B5_\%D1\%82\%D1\%80\%D0\%B5\%D0\%B1\%D0\%BE\%D0\%B2\%D0\%B0\%D0\%BD\%D0\%B8\%D1\%8F.pdf\#nameddest=doc-06-\%D1\%81\%D0\%B2\%D0\%BE\%D0\%B4\%D0\%BD\%D1\%8B\%D0\%B5-\%D0\%BF\%D1\%80\%D0\%BE\%D0\%B3\%D1\%80\%D0\%B0\%D0\%BC\%D0\%BC\%D0\%BD\%D1\%8B\%D0\%B5-\%D1\%82\%D1\%80\%D0\%B5\%D0\%B1\%D0\%BE\%D0\%B2\%D0\%B0\%D0\%BD\%D0\%B8\%D1\%8F}{DOC-06}; \href{../../02_\%D0\%A2\%D0\%B5\%D1\%85\%D0\%BD\%D0\%B8\%D1\%87\%D0\%B5\%D1\%81\%D0\%BA\%D0\%BE\%D0\%B5_\%D0\%B7\%D0\%B0\%D0\%B4\%D0\%B0\%D0\%BD\%D0\%B8\%D0\%B5/\%D0\%A2\%D0\%B5\%D1\%85\%D0\%BD\%D0\%B8\%D1\%87\%D0\%B5\%D1\%81\%D0\%BA\%D0\%BE\%D0\%B5_\%D0\%B7\%D0\%B0\%D0\%B4\%D0\%B0\%D0\%BD\%D0\%B8\%D0\%B5.pdf}{ТЗ, раздел «Архитектура и взаимодействие»}; карточки доменов в \texttt{notes/domains/} \\
Плановый срок & v1.0 --- 15.11.2026 \\
\end{longtable}
}

\section{Назначение и границы}\label{1-ux43dux430ux437ux43dux430ux447ux435ux43dux438ux435-ux438-ux433ux440ux430ux43dux438ux446ux44b}

\hypertarget{1-назначение-и-границы}{}

Документ описывает архитектуру ПО криптотренажёра mephi-12: состав программных компонентов, их размещение по слоям, связи и развёртывание. Он показывает, где в системе размещаются обязательства, определённые \href{../07_\%D0\%9F\%D1\%80\%D0\%BE\%D0\%B3\%D1\%80\%D0\%B0\%D0\%BC\%D0\%BC\%D0\%BD\%D1\%8B\%D0\%B5_\%D1\%82\%D1\%80\%D0\%B5\%D0\%B1\%D0\%BE\%D0\%B2\%D0\%B0\%D0\%BD\%D0\%B8\%D1\%8F/DOC-06_\%D0\%9F\%D1\%80\%D0\%BE\%D0\%B3\%D1\%80\%D0\%B0\%D0\%BC\%D0\%BC\%D0\%BD\%D1\%8B\%D0\%B5_\%D1\%82\%D1\%80\%D0\%B5\%D0\%B1\%D0\%BE\%D0\%B2\%D0\%B0\%D0\%BD\%D0\%B8\%D1\%8F.pdf\#nameddest=doc-06-\%D1\%81\%D0\%B2\%D0\%BE\%D0\%B4\%D0\%BD\%D1\%8B\%D0\%B5-\%D0\%BF\%D1\%80\%D0\%BE\%D0\%B3\%D1\%80\%D0\%B0\%D0\%BC\%D0\%BC\%D0\%BD\%D1\%8B\%D0\%B5-\%D1\%82\%D1\%80\%D0\%B5\%D0\%B1\%D0\%BE\%D0\%B2\%D0\%B0\%D0\%BD\%D0\%B8\%D1\%8F}{DOC-06} и подсистемными требованиями, и не вводит новых требований. Математические алгоритмы, форматы операций и состав требований здесь не дублируются --- на них даются ссылки.

Источником фактов об реализации служат карточки доменов Спринта 1 (\texttt{notes/domains/*.md}) и каталог эндпоинтов (\texttt{notes/\allowbreak{}api/\allowbreak{}endpoints.\allowbreak{}md}), сверенные с кодом \texttt{mephi-12/backend}. Технологии зафиксированы \href{../../02_\%D0\%A2\%D0\%B5\%D1\%85\%D0\%BD\%D0\%B8\%D1\%87\%D0\%B5\%D1\%81\%D0\%BA\%D0\%BE\%D0\%B5_\%D0\%B7\%D0\%B0\%D0\%B4\%D0\%B0\%D0\%BD\%D0\%B8\%D0\%B5/\%D0\%A2\%D0\%B5\%D1\%85\%D0\%BD\%D0\%B8\%D1\%87\%D0\%B5\%D1\%81\%D0\%BA\%D0\%BE\%D0\%B5_\%D0\%B7\%D0\%B0\%D0\%B4\%D0\%B0\%D0\%BD\%D0\%B8\%D0\%B5.pdf}{ТЗ}: Kotlin 1.9.25, Spring Boot 3.4.0 (Web, JPA, Security, Validation), PostgreSQL 17, Docker/docker-compose.

\section{Общий облик системы}\label{2-ux43eux431ux449ux438ux439-ux43eux431ux43bux438ux43a-ux441ux438ux441ux442ux435ux43cux44b}

\hypertarget{2-общий-облик-системы}{}

Система --- монолитное серверное приложение со слоистой структурой и одной внешней зависимостью: реляционной базой данных. Все компоненты выполняются в одном процессе; отдельные сервисы, брокеры сообщений и внешние хранилища не используются. Разделение на слои определяет направление зависимостей: внешний запрос проходит слои сверху вниз, предметные вычисления не обращаются к приёму операций.

{\def\LTcaptype{none} % do not increment counter
\begin{longtable}[]{@{}
  >{\raggedright\arraybackslash}p{(\linewidth - 4\tabcolsep) * \real{0.2275}}
  >{\raggedright\arraybackslash}p{(\linewidth - 4\tabcolsep) * \real{0.3913}}
  >{\raggedright\arraybackslash}p{(\linewidth - 4\tabcolsep) * \real{0.3812}}@{}}
\toprule\noalign{}
\begin{minipage}[b]{\linewidth}\raggedright
Слой
\end{minipage} & \begin{minipage}[b]{\linewidth}\raggedright
Назначение
\end{minipage} & \begin{minipage}[b]{\linewidth}\raggedright
Основные элементы
\end{minipage} \\
\midrule\noalign{}
\endhead
\bottomrule\noalign{}
\endlastfoot
Приём операций & HTTP-приём, аутентификация, авторизация, разбор параметров, единое представление ошибок & Контроллеры, \texttt{JwtFilter}, \texttt{SecurityConfig}, \texttt{ExceptionResolver} \\
Оркестрация & Пользовательские сценарии: coordination доступа, конфигураций, генерации и проверки & \texttt{AuthService}, \texttt{UserService}, \texttt{GroupConfigService}, \texttt{ProblemDecorator}, \texttt{HistoryService} \\
Предметные домены & Генерация условий и математическая проверка решений & Ранцевые решатели и проверяющий, Эль-Гамаль с математической библиотекой, бэкдор-генератор \\
Хранение & Сохранение и чтение данных, согласованность изменений & JPA-сущности, Spring Data-репозитории, конвертеры списковых величин \\
Фоновый контроль & Завершение просроченных задач по расписанию & \texttt{ProblemAutoFailScheduler} \\
\end{longtable}
}

\section{Компоненты и связи}\label{3-ux43aux43eux43cux43fux43eux43dux435ux43dux442ux44b-ux438-ux441ux432ux44fux437ux438}

\hypertarget{3-компоненты-и-связи}{}

Компоненты соответствуют распределению ответственности \href{../07_\%D0\%9F\%D1\%80\%D0\%BE\%D0\%B3\%D1\%80\%D0\%B0\%D0\%BC\%D0\%BC\%D0\%BD\%D1\%8B\%D0\%B5_\%D1\%82\%D1\%80\%D0\%B5\%D0\%B1\%D0\%BE\%D0\%B2\%D0\%B0\%D0\%BD\%D0\%B8\%D1\%8F/DOC-06_\%D0\%9F\%D1\%80\%D0\%BE\%D0\%B3\%D1\%80\%D0\%B0\%D0\%BC\%D0\%BC\%D0\%BD\%D1\%8B\%D0\%B5_\%D1\%82\%D1\%80\%D0\%B5\%D0\%B1\%D0\%BE\%D0\%B2\%D0\%B0\%D0\%BD\%D0\%B8\%D1\%8F.pdf\#nameddest=2-\%D0\%BE\%D1\%82\%D0\%B2\%D0\%B5\%D1\%82\%D1\%81\%D1\%82\%D0\%B2\%D0\%B5\%D0\%BD\%D0\%BD\%D0\%BE\%D1\%81\%D1\%82\%D1\%8C-\%D0\%B8-\%D1\%81\%D0\%B2\%D1\%8F\%D0\%B7\%D0\%B8}{DOC-06}; ниже они дополнены привязкой к реализации и владельцу. Имена классов --- фактические сведения реализации, а не архитектурное ограничение.

{\def\LTcaptype{none} % do not increment counter
\begin{longtable}[]{@{}
  >{\raggedright\arraybackslash}p{(\linewidth - 6\tabcolsep) * \real{0.2047}}
  >{\raggedright\arraybackslash}p{(\linewidth - 6\tabcolsep) * \real{0.3306}}
  >{\raggedright\arraybackslash}p{(\linewidth - 6\tabcolsep) * \real{0.2601}}
  >{\raggedright\arraybackslash}p{(\linewidth - 6\tabcolsep) * \real{0.2047}}@{}}
\toprule\noalign{}
\begin{minipage}[b]{\linewidth}\raggedright
Компонент
\end{minipage} & \begin{minipage}[b]{\linewidth}\raggedright
Ответственность
\end{minipage} & \begin{minipage}[b]{\linewidth}\raggedright
Реализация
\end{minipage} & \begin{minipage}[b]{\linewidth}\raggedright
Владелец
\end{minipage} \\
\midrule\noalign{}
\endhead
\bottomrule\noalign{}
\endlastfoot
Проверка доступа & Аутентификация по JWT, роли, контекст пользователя & \texttt{AuthController}, \texttt{JwtFilter}, \texttt{JwtParser}, \texttt{SecurityConfig} & Кокорев \\
Пользователи и группы & Учётные записи, вступление в группы, история задач & \texttt{UserController}, \texttt{UserService}, \texttt{UserGroupService} & Кокорев \\
Конфигурации и модерация & Доступный состав наборов, признак включения, просмотр сессий & \texttt{GroupConfigService}, \texttt{ModeratorController} & Самбьюи \\
Управление сессиями & Создание набора, допуск и обработка ответа, итог сессии, чтение & \texttt{ProblemController}, \texttt{ProblemDecorator} & Анохин \\
Контроль срока & Фоновое завершение просроченных задач & \texttt{ProblemAutoFailScheduler} & Анохин \\
Ранцевые обработчики & Генерация и проверка ранцевых задач обоих вариантов & \texttt{BackpackProblemsCheckerServiceImpl}, \texttt{BackpackProblemHandler}, семейство решателей & Георгиев \\
Бэкдор-генератор & Учебная демонстрация бэкдора в ранцевой криптосистеме & \texttt{BackdoorController}, \texttt{BackpackBackDoorDegreesProblemSolverServiceImpl} & Минасов \\
Эль-Гамаль и матбиблиотека & Генерация и проверка задач Эль-Гамаля, общая арифметика & \texttt{EIGamalController}, \texttt{ElGamalServiceImpl}, \texttt{utils/mathLibrary} & Портнов \\
Хранение & Модель данных: 8 сущностей и их связи & JPA-сущности, \texttt{*Dao}-репозитории & Самбьюи \\
\end{longtable}
}

\subsection{Последовательность создания набора}\label{ux43fux43eux441ux43bux435ux434ux43eux432ux430ux442ux435ux43bux44cux43dux43eux441ux442ux44c-ux441ux43eux437ux434ux430ux43dux438ux44f-ux43dux430ux431ux43eux440ux430}

\hypertarget{последовательность-создания-набора}{}

Приём операции передаёт в оркестрацию подтверждённого пользователя и имя конфигурации. Оркестрация через компонент конфигураций получает доступный состав набора, затем для каждого квалификатора вызывает генератор предметного домена. Условия, типы и принадлежность задач сохраняются вместе с сессией; подтверждение выдаётся после сохранения полного набора. Порядок закреплён \href{../07_\%D0\%9F\%D1\%80\%D0\%BE\%D0\%B3\%D1\%80\%D0\%B0\%D0\%BC\%D0\%BC\%D0\%BD\%D1\%8B\%D0\%B5_\%D1\%82\%D1\%80\%D0\%B5\%D0\%B1\%D0\%BE\%D0\%B2\%D0\%B0\%D0\%BD\%D0\%B8\%D1\%8F/DOC-06.2_\%D0\%A1\%D0\%B5\%D1\%81\%D1\%81\%D0\%B8\%D0\%B8_\%D0\%B7\%D0\%B0\%D0\%B4\%D0\%B0\%D1\%87_\%D0\%B8_\%D0\%BF\%D0\%BB\%D0\%B0\%D0\%BD\%D0\%B8\%D1\%80\%D0\%BE\%D0\%B2\%D1\%89\%D0\%B8\%D0\%BA.pdf\#nameddest=2-\%D1\%81\%D0\%BE\%D0\%B7\%D0\%B4\%D0\%B0\%D0\%BD\%D0\%B8\%D0\%B5-\%D1\%81\%D0\%B5\%D1\%81\%D1\%81\%D0\%B8\%D0\%B8}{SWR-SES-002--005}.

\subsection{Последовательность обработки ответа}\label{ux43fux43eux441ux43bux435ux434ux43eux432ux430ux442ux435ux43bux44cux43dux43eux441ux442ux44c-ux43eux431ux440ux430ux431ux43eux442ux43aux438-ux43eux442ux432ux435ux442ux430}

\hypertarget{последовательность-обработки-ответа}{}

Приём фиксирует время получения запроса. Управление сессиями проверяет допуск (принадлежность, состояния, срок), затем формат по схеме типа. Допустимый ответ передаётся предметному проверяющему вместе с сохранённым условием; вердикт согласованно сохраняется с состоянием задачи и итогом сессии. Категории результата различаются машинными кодами \href{../07_\%D0\%9F\%D1\%80\%D0\%BE\%D0\%B3\%D1\%80\%D0\%B0\%D0\%BC\%D0\%BC\%D0\%BD\%D1\%8B\%D0\%B5_\%D1\%82\%D1\%80\%D0\%B5\%D0\%B1\%D0\%BE\%D0\%B2\%D0\%B0\%D0\%BD\%D0\%B8\%D1\%8F/DOC-06.2_\%D0\%A1\%D0\%B5\%D1\%81\%D1\%81\%D0\%B8\%D0\%B8_\%D0\%B7\%D0\%B0\%D0\%B4\%D0\%B0\%D1\%87_\%D0\%B8_\%D0\%BF\%D0\%BB\%D0\%B0\%D0\%BD\%D0\%B8\%D1\%80\%D0\%BE\%D0\%B2\%D1\%89\%D0\%B8\%D0\%BA.pdf\#nameddest=\%D0\%BA\%D0\%BE\%D0\%B4\%D1\%8B-\%D0\%B8-\%D0\%BF\%D1\%80\%D0\%B5\%D0\%B4\%D1\%81\%D1\%82\%D0\%B0\%D0\%B2\%D0\%BB\%D0\%B5\%D0\%BD\%D0\%B8\%D0\%B5-\%D0\%BE\%D1\%88\%D0\%B8\%D0\%B1\%D0\%BE\%D0\%BA}{контракта сессий}. Порядок закреплён \href{../07_\%D0\%9F\%D1\%80\%D0\%BE\%D0\%B3\%D1\%80\%D0\%B0\%D0\%BC\%D0\%BC\%D0\%BD\%D1\%8B\%D0\%B5_\%D1\%82\%D1\%80\%D0\%B5\%D0\%B1\%D0\%BE\%D0\%B2\%D0\%B0\%D0\%BD\%D0\%B8\%D1\%8F/DOC-06.2_\%D0\%A1\%D0\%B5\%D1\%81\%D1\%81\%D0\%B8\%D0\%B8_\%D0\%B7\%D0\%B0\%D0\%B4\%D0\%B0\%D1\%87_\%D0\%B8_\%D0\%BF\%D0\%BB\%D0\%B0\%D0\%BD\%D0\%B8\%D1\%80\%D0\%BE\%D0\%B2\%D1\%89\%D0\%B8\%D0\%BA.pdf\#nameddest=3-\%D0\%BE\%D0\%B1\%D1\%80\%D0\%B0\%D0\%B1\%D0\%BE\%D1\%82\%D0\%BA\%D0\%B0-\%D0\%BE\%D1\%82\%D0\%B2\%D0\%B5\%D1\%82\%D0\%B0}{SWR-SES-007--015}.

\subsection{Фоновый контроль и чтение}\label{ux444ux43eux43dux43eux432ux44bux439-ux43aux43eux43dux442ux440ux43eux43bux44c-ux438-ux447ux442ux435ux43dux438ux435}

\hypertarget{фоновый-контроль-и-чтение}{}

Планировщик выполняется в процессе приложения с периодом 60 секунд и приводит просроченные незавершённые задачи к \texttt{FAILED}, учитывая принятые своевременные ответы; агенты внешних очередей отсутствуют. Чтение собственной сессии возвращает сохранённые состояния и результаты без повторной проверки. Обязательства определены \href{../07_\%D0\%9F\%D1\%80\%D0\%BE\%D0\%B3\%D1\%80\%D0\%B0\%D0\%BC\%D0\%BC\%D0\%BD\%D1\%8B\%D0\%B5_\%D1\%82\%D1\%80\%D0\%B5\%D0\%B1\%D0\%BE\%D0\%B2\%D0\%B0\%D0\%BD\%D0\%B8\%D1\%8F/DOC-06.2_\%D0\%A1\%D0\%B5\%D1\%81\%D1\%81\%D0\%B8\%D0\%B8_\%D0\%B7\%D0\%B0\%D0\%B4\%D0\%B0\%D1\%87_\%D0\%B8_\%D0\%BF\%D0\%BB\%D0\%B0\%D0\%BD\%D0\%B8\%D1\%80\%D0\%BE\%D0\%B2\%D1\%89\%D0\%B8\%D0\%BA.pdf\#nameddest=4-\%D0\%BA\%D0\%BE\%D0\%BD\%D1\%82\%D1\%80\%D0\%BE\%D0\%BB\%D1\%8C-\%D1\%81\%D1\%80\%D0\%BE\%D0\%BA\%D0\%B0-\%D0\%B8-\%D0\%B8\%D1\%82\%D0\%BE\%D0\%B3}{SWR-SES-016--019 и SWR-SES-023}.

\section{Размещение и исполнение}\label{4-ux440ux430ux437ux43cux435ux449ux435ux43dux438ux435-ux438-ux438ux441ux43fux43eux43bux43dux435ux43dux438ux435}

\hypertarget{4-размещение-и-исполнение}{}

Все компоненты выполняются в одном процессе приложения; единственный внешний исполнитель --- СУБД. Предметные вычисления выполняются синхронно в рамках обработки запроса; собственные потоки доменов не создаются. Фоновый контроль использует средство планирования платформы и разделяет хранилище с остальными компонентами, поэтому согласование параллельных действий выполняется транзакциями и блокировками данных, а не обменом сообщениями.

Расширение состава задач выполняется подключением нового предметного обработчика под собственным квалификатором: конфигурация называет квалификатор, оркестрация вызывает обработчик, остальные компоненты не изменяются. Требование к такому подключению задано \href{../07_\%D0\%9F\%D1\%80\%D0\%BE\%D0\%B3\%D1\%80\%D0\%B0\%D0\%BC\%D0\%BC\%D0\%BD\%D1\%8B\%D0\%B5_\%D1\%82\%D1\%80\%D0\%B5\%D0\%B1\%D0\%BE\%D0\%B2\%D0\%B0\%D0\%BD\%D0\%B8\%D1\%8F/DOC-06_\%D0\%9F\%D1\%80\%D0\%BE\%D0\%B3\%D1\%80\%D0\%B0\%D0\%BC\%D0\%BC\%D0\%BD\%D1\%8B\%D0\%B5_\%D1\%82\%D1\%80\%D0\%B5\%D0\%B1\%D0\%BE\%D0\%B2\%D0\%B0\%D0\%BD\%D0\%B8\%D1\%8F.pdf\#nameddest=swr-sys-004-\%D0\%BA\%D0\%BE\%D0\%BD\%D1\%82\%D1\%80\%D0\%B0\%D0\%BA\%D1\%82-\%D0\%BF\%D1\%80\%D0\%B5\%D0\%B4\%D0\%BC\%D0\%B5\%D1\%82\%D0\%BD\%D0\%BE\%D0\%B3\%D0\%BE-\%D0\%BE\%D0\%B1\%D1\%80\%D0\%B0\%D0\%B1\%D0\%BE\%D1\%82\%D1\%87\%D0\%B8\%D0\%BA\%D0\%B0}{контрактом предметного обработчика}.

\section{Развёртывание}\label{5-ux440ux430ux437ux432ux451ux440ux442ux44bux432ux430ux43dux438ux435}

\hypertarget{5-развёртывание}{}

Развёртывание выполняется docker-compose из двух служб: контейнер СУБД и контейнер приложения, стартующий после готовности СУБД.

{\def\LTcaptype{none} % do not increment counter
\begin{longtable}[]{@{}
  >{\raggedright\arraybackslash}p{(\linewidth - 2\tabcolsep) * \real{0.3200}}
  >{\raggedright\arraybackslash}p{(\linewidth - 2\tabcolsep) * \real{0.6800}}@{}}
\toprule\noalign{}
\begin{minipage}[b]{\linewidth}\raggedright
Элемент
\end{minipage} & \begin{minipage}[b]{\linewidth}\raggedright
Значение
\end{minipage} \\
\midrule\noalign{}
\endhead
\bottomrule\noalign{}
\endlastfoot
Контейнер СУБД & PostgreSQL 17 (alpine), база \texttt{mephi12}, порт 5432, контроль готовности \texttt{pg\_isready}, том данных \texttt{mephi12-data} \\
Контейнер приложения & Сборка Gradle на JDK 17, исполнение на OpenJDK 17, порт 8083, базовый путь \texttt{/api} \\
Параметры окружения & \texttt{JWT\_SECRET}, \texttt{ACCESS\_\allowbreak{}TOKEN\_\allowbreak{}LIFETIME}, \texttt{DATABASE\_URL}, \texttt{DATABASE\_USER}, \texttt{DATABASE\_PASSWORD} \\
Схема базы & Создаётся средствами JPA по описанию сущностей \\
\end{longtable}
}

Обязательства развёртывания заданы {[}ТЗ, \texttt{TS.REQ.REL.DEPLOY}{]}: старт приложения только после готовности СУБД и сохранность данных между перезапусками томом данных. Секреты передаются только переменными окружения и не включаются в код или репозиторий ({[}\texttt{TS.REQ.SEC.SECRETS}{]}).

\section{Соответствие компонентов требованиям}\label{6-ux441ux43eux43eux442ux432ux435ux442ux441ux442ux432ux438ux435-ux43aux43eux43cux43fux43eux43dux435ux43dux442ux43eux432-ux442ux440ux435ux431ux43eux432ux430ux43dux438ux44fux43c}

\hypertarget{6-соответствие-компонентов-требованиям}{}

Таблица связывает компоненты с документами требований, обязательства которых они размещают. Полная матрица SR → SWR → TC ведётся в DOC-09.

{\def\LTcaptype{none} % do not increment counter
\begin{longtable}[]{@{}
  >{\raggedright\arraybackslash}p{(\linewidth - 4\tabcolsep) * \real{0.3698}}
  >{\raggedright\arraybackslash}p{(\linewidth - 4\tabcolsep) * \real{0.3156}}
  >{\raggedright\arraybackslash}p{(\linewidth - 4\tabcolsep) * \real{0.3146}}@{}}
\toprule\noalign{}
\begin{minipage}[b]{\linewidth}\raggedright
Компонент
\end{minipage} & \begin{minipage}[b]{\linewidth}\raggedright
Требования
\end{minipage} & \begin{minipage}[b]{\linewidth}\raggedright
Документ
\end{minipage} \\
\midrule\noalign{}
\endhead
\bottomrule\noalign{}
\endlastfoot
Проверка доступа; пользователи и группы & \texttt{SR.AUTH.*}, \texttt{SWR.AUTH.*} & DOC-06.1 (Кокорев) \\
Управление сессиями; контроль срока & \texttt{SR-SES-*}, \texttt{SWR-SES-*} & \href{../07_\%D0\%9F\%D1\%80\%D0\%BE\%D0\%B3\%D1\%80\%D0\%B0\%D0\%BC\%D0\%BC\%D0\%BD\%D1\%8B\%D0\%B5_\%D1\%82\%D1\%80\%D0\%B5\%D0\%B1\%D0\%BE\%D0\%B2\%D0\%B0\%D0\%BD\%D0\%B8\%D1\%8F/DOC-06.2_\%D0\%A1\%D0\%B5\%D1\%81\%D1\%81\%D0\%B8\%D0\%B8_\%D0\%B7\%D0\%B0\%D0\%B4\%D0\%B0\%D1\%87_\%D0\%B8_\%D0\%BF\%D0\%BB\%D0\%B0\%D0\%BD\%D0\%B8\%D1\%80\%D0\%BE\%D0\%B2\%D1\%89\%D0\%B8\%D0\%BA.pdf\#nameddest=doc-062-\%D1\%81\%D0\%B5\%D1\%81\%D1\%81\%D0\%B8\%D0\%B8-\%D0\%B7\%D0\%B0\%D0\%B4\%D0\%B0\%D1\%87-\%D0\%B8-\%D0\%BF\%D0\%BB\%D0\%B0\%D0\%BD\%D0\%B8\%D1\%80\%D0\%BE\%D0\%B2\%D1\%89\%D0\%B8\%D0\%BA}{DOC-06.2} (Анохин) \\
Ранцевые обработчики & \texttt{SR-KNAP-*}, \texttt{SWR-KNAP-*} & \href{../07_\%D0\%9F\%D1\%80\%D0\%BE\%D0\%B3\%D1\%80\%D0\%B0\%D0\%BC\%D0\%BC\%D0\%BD\%D1\%8B\%D0\%B5_\%D1\%82\%D1\%80\%D0\%B5\%D0\%B1\%D0\%BE\%D0\%B2\%D0\%B0\%D0\%BD\%D0\%B8\%D1\%8F/DOC-06.3_\%D0\%A0\%D0\%B0\%D0\%BD\%D1\%86\%D0\%B5\%D0\%B2\%D1\%8B\%D0\%B5_\%D0\%BA\%D1\%80\%D0\%B8\%D0\%BF\%D1\%82\%D0\%BE\%D1\%81\%D0\%B8\%D1\%81\%D1\%82\%D0\%B5\%D0\%BC\%D1\%8B.pdf\#nameddest=doc-063-\%D1\%80\%D0\%B0\%D0\%BD\%D1\%86\%D0\%B5\%D0\%B2\%D1\%8B\%D0\%B5-\%D0\%BA\%D1\%80\%D0\%B8\%D0\%BF\%D1\%82\%D0\%BE\%D1\%81\%D0\%B8\%D1\%81\%D1\%82\%D0\%B5\%D0\%BC\%D1\%8B}{DOC-06.3} (Георгиев) \\
Бэкдор-генератор & \texttt{SR.BDR.*}, \texttt{SWR.BDR.*} & DOC-06.4 (Минасов) \\
Эль-Гамаль и матбиблиотека & \texttt{SR.ELG.*}, \texttt{SWR.ELG.*} & DOC-06.5 (Портнов) \\
Конфигурации, модерация, хранение & \texttt{SR.CFG.*}, \texttt{SWR.CFG.*} & DOC-06.6 (Самбьюи) \\
Общие свойства взаимодействия & \texttt{SR-SYS-*}, \texttt{SWR-SYS-*} & \href{../07_\%D0\%9F\%D1\%80\%D0\%BE\%D0\%B3\%D1\%80\%D0\%B0\%D0\%BC\%D0\%BC\%D0\%BD\%D1\%8B\%D0\%B5_\%D1\%82\%D1\%80\%D0\%B5\%D0\%B1\%D0\%BE\%D0\%B2\%D0\%B0\%D0\%BD\%D0\%B8\%D1\%8F/DOC-06_\%D0\%9F\%D1\%80\%D0\%BE\%D0\%B3\%D1\%80\%D0\%B0\%D0\%BC\%D0\%BC\%D0\%BD\%D1\%8B\%D0\%B5_\%D1\%82\%D1\%80\%D0\%B5\%D0\%B1\%D0\%BE\%D0\%B2\%D0\%B0\%D0\%BD\%D0\%B8\%D1\%8F.pdf\#nameddest=doc-06-\%D1\%81\%D0\%B2\%D0\%BE\%D0\%B4\%D0\%BD\%D1\%8B\%D0\%B5-\%D0\%BF\%D1\%80\%D0\%BE\%D0\%B3\%D1\%80\%D0\%B0\%D0\%BC\%D0\%BC\%D0\%BD\%D1\%8B\%D0\%B5-\%D1\%82\%D1\%80\%D0\%B5\%D0\%B1\%D0\%BE\%D0\%B2\%D0\%B0\%D0\%BD\%D0\%B8\%D1\%8F}{DOC-06} \\
\end{longtable}
}

\section{Известные ограничения реализации}\label{7-ux438ux437ux432ux435ux441ux442ux43dux44bux435-ux43eux433ux440ux430ux43dux438ux447ux435ux43dux438ux44f-ux440ux435ux430ux43bux438ux437ux430ux446ux438ux438}

\hypertarget{7-известные-ограничения-реализации}{}

Ограничения зафиксированы \href{../../02_\%D0\%A2\%D0\%B5\%D1\%85\%D0\%BD\%D0\%B8\%D1\%87\%D0\%B5\%D1\%81\%D0\%BA\%D0\%BE\%D0\%B5_\%D0\%B7\%D0\%B0\%D0\%B4\%D0\%B0\%D0\%BD\%D0\%B8\%D0\%B5/\%D0\%A2\%D0\%B5\%D1\%85\%D0\%BD\%D0\%B8\%D1\%87\%D0\%B5\%D1\%81\%D0\%BA\%D0\%BE\%D0\%B5_\%D0\%B7\%D0\%B0\%D0\%B4\%D0\%B0\%D0\%BD\%D0\%B8\%D0\%B5.pdf}{ТЗ, п. 2.3} и карточками доменов; они не меняют описанную архитектуру, а помечают различия между ней и текущим кодом.

\begin{enumerate}
\def\labelenumi{\arabic{enumi}.}
\tightlist
\item
  Схема базы создаётся режимом обновления JPA-описания; инструменты миграции не применяются.
\item
  Правило доступа к функции модерации фактически не отличает роли: эндпоинт просмотра сессий доступен любой роли \texttt{USER}, а выдача роли \texttt{ADMIN} через API невозможна.
\item
  Настройка CORS разрешает все источники с передачей учётных данных; область применения системы --- учебная.
\item
  Управление группами и конфигурациями на уровне API не покрывает создание, изменение и удаление --- выполняется средствами СУБД.
\item
  Средство наблюдения (Actuator) входит в зависимости, но не настроено; собственные метрики и зонды готовности приложения не экспортируются.
\end{enumerate}

История: 07.11.2026 --- подготовлен проект 0.5: слои, компоненты и связи, последовательности операций, развёртывание и соответствие требованиям.


\end{document}
